\section{Trace study}
\label{sec:trace-method}

\subsection{Pinned public source}

The source is the complete \texttt{tau2\_airline}, \texttt{tau2\_retail}, and
\texttt{tau2\_telecom} subset of the public Exgentic agent-trace dataset. The
dataset revision is
\sha{70036b93a04e61b0ea2706a68b962f4f26774587}
and the Parquet conversion revision is
\sha{f7c94012d0bfbf66fe4d6ed627699508bbb555ff}; all 19 retained shard hashes
match commit-resolved URLs \cite{exgentic2026traces}. The domain labels correspond to the airline, retail,
and telecom environments released with $\tau^2$-Bench
\cite{barres2026tau2bench}, which builds on the original $\tau$-bench framework
\cite{yao2024taubench}. The selected panel contains
\TraceSessions{} sessions and \TraceSpans{} recorded LLM spans across four
harnesses.

\claimid{RC-015}Each span completion becomes one candidate control event. Its route key is
derived from the recorded outcome: final or text, error, one named tool, or a
generic multi-tool outcome. It omits benchmark and harness semantics,
state-machine node, schema and version, arguments, policy context, and the tool
identities inside a multi-tool outcome. The extraction retains timestamps,
public identifiers, counts, lengths, and route labels. It excludes prompt text,
tool arguments, and tool results. Source Parquet files and both derived outputs
are content-hashed.

The selected data contain 70 route labels. Eight labels occur in more than one
benchmark domain, and 28 occur under more than one harness. The generic
\texttt{tool:<multi>} label alone covers 325 spans and 122 distinct recorded
tool-name encodings. The source audit also retains 52 failed-status spans, 52
nonpositive-duration spans, and 617 overlapping span starts. We do not remove
these observations after inspecting outcomes.

This construction treats a model completion as the point at which a control
transition could become ready. The route key is a declared conditioning proxy,
not proof of executable compatibility. The construction does not claim that
the recorded harnesses implemented our transition or that their span durations
equal model service times.

\subsection{Stationary replay}

We use the prospectively frozen stationary swarm model from trace replay 003. Session
arrivals follow a homogeneous Poisson process. A session template is sampled
uniformly from the fixed empirical panel. For target mean active population
$C$, the arrival rate is $C$ divided by mean template duration. Arrivals begin
one maximum template duration before measurement so that the retained
60-second interval is in steady state under the model. Fixed partitions are
origin-aligned half-open windows of width $\delta$; $F$ is conditional on that
phase, whereas $\pstar$ is not.

The exact-packing grid contains:

\begin{itemize}
  \item $C\in\{1{,}000,10{,}000,100{,}000\}$;
  \item $\delta\in\{10,25,50,100,250\}$ ms;
  \item pooled, event-class, and route-key grouping;
  \item $K\in\{32,64,128,256\}$; and
  \item three replay seeds from root seed 20260811.
\end{itemize}

\claimid{RC-007}Every generated event is retained. The three population values crossed with
three seeds produce nine generated swarms. Reusing each swarm across deadlines,
groupings, and thresholds yields \TraceReplayCells{} design cells and
\TraceReplayRepetitions{} cell-seed rows. For each row we compute $F$, $\pstar$,
$U$, the exact batch count, and

\begin{equation}
  G=\frac{\pstar-F}{U-F}\quad\text{when }U>F.
  \label{eq:gap-closure}
\end{equation}

$G$ describes how much of the local boundary-alignment opportunity is jointly
packable. The reported primary $G$ is the mean of the three per-seed ratios,
not the ratio formed from three mean shares.

\subsection{Prospective freeze and validity gates}

The exact solver, integer-clock contract, inputs, grid, primary cell, and
directional hypothesis were frozen before any exact trace outcome was
computed according to the artifact-recorded chronology. The plan was not
externally registered or independently timestamped. Earlier fixed-window and local-bound results were permitted design
data. Five invariants gate interpretation: $F\leq\pstar\leq U$; monotonicity
under grouping coarsening; monotonicity with $\delta$; antitonicity with
$K$; and equality of all three shares whenever the previously computed lower
and upper bounds coincide.

The primary cell was $C=100{,}000$, route-key grouping, $K=256$, and
$\delta=50$ ms.
The directional pilot hypothesis was only $\pstar>F$. The three seeds quantify
Monte Carlo variation under one fixed panel and arrival model. They are not
independent traces from a deployment population, so no population $p$-value or
confidence interval is reported.

\begin{table}[t]
  \centering
  \footnotesize
  \begin{tabularx}{\columnwidth}{@{}lX@{}}
    \toprule
    Frozen item & Value \\
    \midrule
    Source & Exgentic tau2 panel, two commit hashes \\
    Empirical unit & one of \TraceSessions{} session templates \\
    Event & recorded LLM-span completion \\
    Grouping & outcome-derived route-key proxy \\
    Replay & stationary Poisson session arrivals \\
    Clock & integer nanoseconds, inclusive endpoints \\
    Capacity/service & unlimited / zero \\
    Uncertainty & three seeds conditional on one panel \\
    \bottomrule
  \end{tabularx}
  \caption{Trace-study contract. The replay is a controlled load model rather
  than a model of production arrival statistics.}
  \label{tab:trace-contract}
\end{table}
